\section*{Chapitre \ref{ch:m3:crypto-data-and-infrastructure} --- Données crypto et infrastructure}

\begin{solution}{exo:m3:crypto-data-and-infrastructure:1}
On ignore $(95, 99)$ (déjà dans le carnet)~; on applique $(101, 103)$, l'identifiant du carnet devenant 103~; $(105, 106)$ commence après 104, donc la mise à jour 104 a été manquée~:
on abandonne le carnet et on resynchronise à partir d'un nouvel instantané.
\end{solution}

\begin{solution}{exo:m3:crypto-data-and-infrastructure:2}
Lorsqu'il faut interroger rapidement et de façon fiable l'état historique à des blocs passés quelconques (soldes, positions, stockage des contrats à ce moment-là), pour la recherche,
la réconciliation ou le traçage~; un \omterm{def:m3:crypto-data-and-infrastructure:node}{nœud complet} peut en reconstituer une partie, mais lentement.
\end{solution}

\begin{solution}{exo:m3:crypto-data-and-infrastructure:3}
«~10014~» «~10026~» «~9995~» «~9987~» concaténés, les cours vendeurs d'abord du plus bas au plus haut, puis les cours acheteurs du plus haut au plus bas~; leur CRC32 vaut 3\,348\,617\,501.
\end{solution}

\begin{solution}{exo:m3:crypto-data-and-infrastructure:4}
Environ $pK = 0.1\% \times 20 = 2\%$ du temps (1.9~\% dans la simulation)~; une erreur typique dure environ 20 messages (21 simulés), jusqu'à
la prochaine mise à jour du niveau perdu.
\end{solution}

\begin{solution}{exo:m3:crypto-data-and-infrastructure:5}
Une stratégie ne peut agir que sur ce qu'elle a reçu~; les horodatages de la bourse indiquent quand la plateforme signale qu'un événement s'est produit, plus tôt que la société ne peut le savoir,
si bien que les décisions datées par ces horodatages regardent dans le futur, d'une durée égale à la latence du flux.
\end{solution}

\begin{solution}{exo:m3:crypto-data-and-infrastructure:6}
On le détecte à partir des hachages de blocs, on supprime ou on marque comme orphelins les enregistrements des deux blocs abandonnés, on relit les nouveaux blocs à ces hauteurs, et
on propage les corrections à tout ce qui a été calculé à partir d'eux~; les enregistrements récents doivent porter leur hachage de bloc et un indicateur de \omterm{def:m3:blockchains-for-traders:finality}{finalité}.
\end{solution}

\begin{solution}{exo:m3:crypto-data-and-infrastructure:7}
Environ 20 messages~: la part des données périmées vaut $pR$ et la part des erreurs silencieuses $pK$, égales lorsque $R = K = 20$ (toutes deux à 1.92~\% dans la simulation).
\end{solution}

\begin{solution}{exo:m3:crypto-data-and-infrastructure:8}
Enregistrer ne revient pas à reconstruire~: l'enregistreur du fournisseur peut perdre des messages, et son jeu de données reflète la façon dont il a traité les lacunes~; il faut vérifier les carnets croisés et les ruptures de
séquence, les trous et les resynchronisations, puis comparer avec une seconde source.
\end{solution}

\begin{solution}{pb:m3:crypto-data-and-infrastructure:1}
\textbf{1.} Environ 1.9~\% des messages. \textbf{2.} Environ 21 messages. \textbf{3.} Chaque message met à jour le niveau perdu avec une probabilité $1/20$, de sorte que l'attente suit
une loi géométrique de moyenne environ 20. \textbf{4.} Coter autour d'un milieu erroné, croire qu'un niveau existe alors qu'il n'existe pas, ou manquer un arbitrage qui n'existe pas~; dans le
pire cas, négocier contre un carnet croisé fantôme. \textbf{5.} Une suppression perdue laisse un niveau fantôme qui subsiste jusqu'à ce que ce prix précis soit de nouveau mis à jour,
ce qui, si le cours s'éloigne, peut ne jamais se produire. \textbf{6.} Environ 4.75~\%. \textbf{7.} Resynchroniser (50 messages) prend plus de temps qu'une erreur silencieuse typique
(20)~; le client prudent troque la connaissance de son état périmé contre un état périmé plus longtemps. \textbf{8.} Il retire ou élargit ses cotations et marque ses données comme manquantes.
\textbf{9.} En appliquant les événements mis en tampon sur le nouvel instantané plutôt qu'en attendant de nouveaux événements, la fenêtre de péremption se réduit à la latence de l'instantané.
\textbf{10.} Protection contre les bogues propres au client et contre les erreurs d'identifiants, et vérification du carnet qu'il calcule. \textbf{11.} Vérifier les identifiants de
mise à jour pour repérer les trous, rechercher les carnets croisés ou verrouillés, comparer le haut du carnet avec les transactions, et comparer avec une seconde source. \textbf{12.} Une mise à jour perdue,
le plus souvent une suppression manquée ou un instantané tardif, puisqu'un carnet d'une plateforme réelle ne peut pas se croiser. \textbf{13.} L'horodatage de réception. \textbf{14.} En s'appuyant sur leurs horodatages
de bourse corrigés de la latence mesurée de chaque flux, ou sur les horodatages de réception en un seul emplacement. \textbf{15.} Le moment où chaque resynchronisation a commencé et s'est achevée,
afin que les utilisateurs puissent écarter ou signaler ces périodes. \textbf{16.} Oui~: un carnet périmé sait qu'il ne peut pas être fiable~; un carnet erroné, lui, négocie. \textbf{17.} Un fournisseur suffit
pour des requêtes ponctuelles et la recherche~; le trading sensible à la latence, les forts volumes de requêtes et l'indépendance exigent ses propres nœuds. \textbf{18.} Plus difficile~: de nombreuses plateformes, chacune avec
ses propres conventions de flux, des infrastructures cloud publiques, pas de flux consolidé, des chaînes qui se réorganisent~; plus facile~: des données gratuites, publiques et en profondeur complète de chaque plateforme et de chaque chaîne.
\textbf{19.} \emph{Résultat chiffré~:} pour un taux de perte de 0.1~\% et 20 niveaux, 1.9~\% du temps silencieusement faux sans contrôles (erreurs d'une durée d'environ 21 messages),
contre 4.75~\% de données sciemment périmées, la resynchronisation prenant 50 messages. \textbf{20.} La copie que la société tient du carnet d'une plateforme, aussi bonne que sa gestion des
messages qu'elle n'a pas reçus.
\end{solution}

\begin{solution}{iq:m3:crypto-data-and-infrastructure:1}
Mettre le flux en tampon, récupérer un instantané, écarter les événements mis en tampon déjà contenus dans celui-ci, appliquer les autres dans l'ordre en vérifiant que chaque événement commence juste après
l'identifiant du carnet, fixer les niveaux en quantités absolues (zéro supprime), et resynchroniser à la moindre rupture~; vérifier à l'aide des sommes de contrôle lorsque la plateforme les publie.
\iqlookfor{la règle de séquence et la resynchronisation.}
\end{solution}

\begin{solution}{iq:m3:crypto-data-and-infrastructure:2}
L'horodatage de la bourse est l'heure de la plateforme pour l'événement (ou le message)~; l'horodatage de réception est le moment où la société l'a reçu. Utiliser les horodatages de réception pour ce qu'une stratégie savait
et les horodatages de bourse pour aligner les plateformes et mesurer la latence.
\iqlookfor{connaissance contre alignement.}
\end{solution}

\begin{solution}{iq:m3:crypto-data-and-infrastructure:3}
Lire les blocs depuis ses propres nœuds~; décoder les journaux de swaps par adresse de pool et signature~; stocker avec le numéro de bloc, le hachage, l'indice du journal et la \omterm{def:m3:blockchains-for-traders:finality}{finalité}~; gérer les réorganisations en
comparant les hachages parents et réécrivant les plages orphelines~; compléter l'historique à partir d'un \omterm{def:m3:crypto-data-and-infrastructure:node}{nœud d'archive}~; associer des versions de décodeurs aux mises à niveau des contrats.
\iqlookfor{gestion des réorganisations et rigueur du décodage.}
\end{solution}

\begin{solution}{iq:m3:crypto-data-and-infrastructure:4}
Journaliser le carnet et le message à chaque échec~; vérifier la règle de formatage (nombre de décimales, zéros de fin, nombre de niveaux), l'ordre des niveaux, le traitement des
quantités nulles et des niveaux au-delà de la profondeur de l'instantané, et les pertes ou réordonnancements de messages autour des échecs~; reproduire hors ligne à partir du flux enregistré.
\iqlookfor{formatage d'abord, puis séquence.}
\end{solution}

\begin{solution}{iq:m3:crypto-data-and-infrastructure:5}
Comparer les horodatages de bourse et de réception sur les messages des périodes calmes, mesurer les allers-retours de requêtes peu coûteuses depuis des machines situées dans plusieurs régions, et observer la distribution plutôt
que la moyenne, surtout aux heures chargées~; synchroniser d'abord les horloges de la société.
\iqlookfor{rigueur des horloges et distributions.}
\end{solution}

\begin{solution}{iq:m3:crypto-data-and-infrastructure:6}
Des adaptateurs par plateforme qui normalisent les messages et conservent les deux horodatages~; des constructeurs de carnets avec contrôles de séquence et de somme de contrôle~; des grappes de nœuds et des \omterm{def:m3:crypto-data-and-infrastructure:log}{indexeurs} par chaîne~; un
bus normalisé vers les stratégies avec des indicateurs de péremption~; enregistrement des messages bruts pour le rejeu~; surveillance des trous, de la latence et des resynchronisations~; et déploiement dans les régions
les plus proches de chaque plateforme.
\iqlookfor{normaliser, vérifier, signaler et enregistrer.}
\end{solution}
